\begin{helper}
Fix a history cell \(H\), a finite terminal coordinate set \(U\), one blue
support label \(B\), and one red support label \(J\).  Let \(\mathcal T\) be
the finite branch-label type, let \(\mathcal P\) be the finite transcript
family, and write
\[
\tau=(P_\tau,A_\tau,Z_\tau).
\]
For each branch label \(\beta\), let
\(\operatorname{blueTrace}(\beta)\subseteq U\) be the complete terminal blue
truth table, let \(R_{\beta,J}\subseteq U\) be the red support decoded from
\(J\) after the branch is fixed, and let
\[
O_\beta\subseteq U
\]
be the branch-indexed set of terminal coordinates that are still
staged-open candidates after the fixed blue trace has been exposed.  Let
\(\mathsf{Realized}(\beta,\tau)\) mean that the fixed \(B,J\)
branch/transcript cell is nonempty.

Assume the complete terminal product law over \(H\): each complete terminal
assignment \(A\subseteq U\) has conditional mass
\[
p^{|A|}(1-p)^{|U|-|A|}.
\]
Assume \(0\le p\le1/2\).  Assume \(B\) is a blue \(u_B\)-subset of \(U\).
For every realized pair \((\beta,\tau)\), assume
\[
|R_{\beta,J}|=u_R,\qquad R_{\beta,J}\subseteq U,
\]
that \(R_{\beta,J}\) and \(Z_\tau\) are red-only, and that
\[
P_\tau,A_\tau\subseteq U,\qquad P_\tau\cap A_\tau=\emptyset,
\]
\[
B\cup R_{\beta,J}\subseteq P_\tau,\qquad
Z_\tau\subseteq A_\tau.
\]
Assume also that on blue terminal coordinates, \(P_\tau\) and \(A_\tau\)
record the complete branch trace:
\[
e\in\operatorname{blueTrace}(\beta)\Longleftrightarrow e\in P_\tau,
\qquad
e\notin\operatorname{blueTrace}(\beta)\Longleftrightarrow e\in A_\tau .
\]
Assume branch labels are functional from this complete terminal blue trace:
if two labels have the same truth table on every blue coordinate of \(U\),
then the two labels are equal.  Consequently, by
\(\noderef{FixedSetHistoryCellRedSupportDeterminedByBlueTrace}\), for fixed
\(J\) the decoded red support \(R_{\beta,J}\) is determined by the same
blue-trace data.

Assume the fixed-support exhaustion facts for the selected absence scan.  For
every realized \((\beta,\tau)\),
\[
O_\beta\subseteq U,\qquad
B\cup R_{\beta,J}\subseteq O_\beta,\qquad
Z_\tau\subseteq O_\beta,
\]
and every open candidate forced present by the cell is one of the fixed
support coordinates:
\[
e\in O_\beta,\ e\in P_\tau
  \quad\Longrightarrow\quad
e\in B\cup R_{\beta,J}.
\]
For each selected split \(e\in Z_\tau\), let
\(\mathsf{Sibling}(A,\beta,\tau,e)\) mean that \(A\subseteq U\) is a
complete terminal assignment in the corresponding selected present sibling.
To formalize the mass conservation of the residual tree without explicitly
constructing its exposure order, we assume a coverage and partition hypothesis.
For any realized cell \((\beta,\tau)\), let its \emph{enlarged undivided cylinder} be
the set of complete assignments \(A\subseteq U\) that match the blue trace
\(\operatorname{blueTrace}(\beta)\) on all blue coordinates outside \(Z_\tau\), contain all forced present coordinates
\(P_\tau\), and are disjoint from \(A_\tau\setminus Z_\tau\).
Assume that this enlarged undivided cylinder is exactly partitioned by the union of the
compatible assignments for \((\beta,\tau)\) and the selected present siblings
for all \(e\in Z_\tau\).
Furthermore, assume that for any two distinct realized cells \((\beta,\tau)\)
and \((\beta',\tau')\), their enlarged undivided cylinders are disjoint.
Because the present sibling assignments are pruned, we also assume that
if \(\mathsf{Sibling}(A,\beta,\tau,e)\), then \(A\) is not compatible with
any realized cell \((\beta',\tau')\).

Finally assume that compatibility with a complete terminal assignment is
exactly the cylinder condition for \(P_\tau,A_\tau\) together with this blue
truth table, and that the branch choice and prefix scan are functional from
the assignment: if \(A\subseteq U\) is compatible with a realized cell, then
the chosen branch is \(\beta\) and the scan output is \(\tau\).

Then, for this fixed pair \(B,J\), there is a single finite residual leaf
type \(\Lambda\), a nonnegative leaf mass \(\nu:\Lambda\to\mathbb R\), and
an assignment
\[
\iota(\beta,\tau)\in\Lambda
\]
for every formal branch/transcript pair such that
\[
\sum_{\lambda\in\Lambda}\nu(\lambda)\le1,
\]
the assignment is injective on realized cells, and for every realized
\((\beta,\tau)\),
\[
p^{|P_\tau|}(1-p)^{|A_\tau|}
\le
p^{u_R}p^{u_B}
(1-p)^{|Z_\tau|}\,\nu(\iota(\beta,\tau)).
\]
All absent blue-coordinate factors from the complete branch trace remain
inside the residual leaf mass \(\nu\); they are not divided out
branch-by-branch.
\end{helper}
\begin{proof}
Put
\[
q=1-p.
\]
From \(0\le p\le1/2\) we have \(1/2\le q\le1\).  The construction is made
once for the fixed support labels \(B,J\).  The residual tree is therefore a
single tree for the whole fixed-\((B,J)\) summation, not a separate tree for
each realized cell.

First note what the coverage and partition hypothesis implies about the
probability mass.  For a realized cell \((\beta,\tau)\), its enlarged
undivided cylinder \(C'_{\beta,\tau}\) requires the presence of \(P_\tau\)
and the absence of \(A_\tau\setminus Z_\tau\), but places no constraint on
the selected residual coordinates \(Z_\tau\).  It still fixes the complete
blue trace of \(\beta\): since \(Z_\tau\) is red-only, releasing \(Z_\tau\)
does not release any blue absent coordinate.  The probability mass of this
cylinder under the complete terminal product law is
\[
P(C'_{\beta,\tau})
 = p^{|P_\tau|}(1-p)^{|A_\tau|-|Z_\tau|}
 = p^{|P_\tau|}q^{|A_\tau|-|Z_\tau|}.
\]
The partition hypothesis states that \(C'_{\beta,\tau}\) is exactly the
disjoint union of the compatible assignments for \((\beta,\tau)\) and the
pruned present siblings \(\mathsf{Sibling}(A,\beta,\tau,e)\) for
\(e\in Z_\tau\).  Furthermore, the enlarged undivided cylinders
\(C'_{\beta,\tau}\) are mutually disjoint across all realized \(\beta\) and
\(\tau\).  For a fixed \(\beta\), the deterministic prefix scan gives
different released cylinders for different transcript outputs.  Across
different branches, the released cylinders are disjoint because they require
different complete blue traces, and releasing \(Z_\tau\) has not released any
blue coordinate.

Let \(E=\bigcup_{\beta,\tau}C'_{\beta,\tau}\) be the disjoint union of these
enlarged undivided cylinders.  Every assignment in \(C'_{\beta,\tau}\) has
its complete blue trace matching \(\beta\), forces the blue support \(B\)
present, and forces the decoded red support \(R_{\beta,J}\) present.
Therefore
\[
E \subseteq \bigcup_{\beta} \Bigl( \{\text{blue trace is } \beta\} \cap \{B \text{ is present}\} \cap \{R_{\beta,J} \text{ is present}\} \Bigr).
\]
Because the decoded red support \(R_{\beta,J}\) depends only on the blue
trace, once the branch trace is fixed the event that \(R_{\beta,J}\) is
present involves only red terminal coordinates and has probability
\(p^{u_R}\).  It is independent of the blue trace and of the event that the
blue support \(B\) is present.  Since the branch blue traces are mutually
exclusive, the total probability is bounded by
\[
P(E)
\le \sum_{\beta}
  P(\text{blue trace is } \beta \text{ and } B \text{ is present})\,p^{u_R}
= P(B \text{ is present})\,p^{u_R}
= p^{u_B}p^{u_R}.
\]

A residual leaf is defined algebraically as one realized terminal block of
this process, together with a distinguished dummy leaf used for non-realized
formal pairs.  The block records the branch label, the transcript output, the
decoded red support, the selected residual absence set, and the remaining
forced present and forced absent terminal values after only the factors from
\(B\), \(R_{\beta,J}\), and \(Z_\tau\) have been removed.  In particular, the
absent blue coordinates in the complete trace are still part of the residual
block.  Because \(\mathcal T\) and \(\mathcal P\) are finite, the set of
realized residual blocks, with the dummy leaf adjoined, is finite.  Let
\(\Lambda\) be this finite leaf set.

Define \(\nu(\lambda)\) for a realized block
\(\lambda=\iota(\beta,\tau)\) by
\[
\nu(\iota(\beta,\tau)) = \begin{cases} 0 & \text{if } p=0 \text{ and } u_R+u_B>0, \\
p^{|P_\tau|-u_R-u_B} q^{|A_\tau|-|Z_\tau|} & \text{otherwise}. \end{cases}
\]
Since \(Z_\tau\subseteq A_\tau\), the exponent
\(|A_\tau|-|Z_\tau|\) is nonnegative.  For the dummy leaf, set the mass to
\(0\).  Nonnegativity is immediate from \(0\le p\) and \(0\le q\).
Except in the trivial \(p=0\) zero-weight case,
\[
\nu(\iota(\beta,\tau))
= p^{-u_R-u_B} P(C'_{\beta,\tau}).
\]
The complete-blue absent factors are included in \(P(C'_{\beta,\tau})\) and
therefore in \(\nu\); no factor \(q^{-|A_{\beta,\text{blue}}|}\) is ever
introduced.  Summing over the disjoint released cylinders gives
\[
\sum_{\beta,\tau}\nu(\iota(\beta,\tau))
= p^{-u_R-u_B}P(E)
\le 1.
\]

We next define \(\iota\).  If \((\beta,\tau)\) is realized, let
\(\iota(\beta,\tau)\) be the residual block determined by that cell.  This
is well-defined even when the cell contains many compatible complete
assignments: the forward compatibility clause says that every compatible
assignment contains \(P_\tau\), avoids \(A_\tau\), and has the blue truth
table \(\operatorname{blueTrace}(\beta)\); the converse clause says that
these cylinder conditions are exactly the compatible assignments for the
realized cell.  Hence the residual block depends only on the data
\((\beta,\tau)\), not on a chosen assignment inside the cylinder.  If
\((\beta,\tau)\) is not realized, set \(\iota(\beta,\tau)\) to the dummy
leaf.

The map \(\iota\) is injective on realized cells.  Suppose
\[
\iota(\beta,\tau)=\iota(\beta',\tau')
\]
with both cells realized.  Equality of residual blocks gives the same
remaining forced terminal values and also records the removed data
\(B\), \(R_{\beta,J}\), \(Z_\tau\) and
\(R_{\beta',J}\), \(Z_{\tau'}\).  Reinsert the fixed present support values
and the selected absent values into the common residual block.  This
produces a complete terminal assignment compatible with both realized cells.
By the functional branch/scan hypothesis, equivalently by
\(\noderef{FixedSetHistoryCellRedFunctionalScanDisjointness}\), the same
complete assignment has one branch label and one scan output.  Therefore
\[
\beta=\beta'\qquad\text{and}\qquad \tau=\tau'.
\]

It remains to check the displayed domination.  Fix a realized cell
\((\beta,\tau)\).  The cell geometry gives
\[
B\cup R_{\beta,J}\subseteq P_\tau,\qquad
Z_\tau\subseteq A_\tau,\qquad P_\tau\cap A_\tau=\emptyset,
\]
\[
|B|=u_B,\qquad |R_{\beta,J}|=u_R.
\]
The blue support \(B\) and red support \(R_{\beta,J}\) are disjoint by their
colors.  The selected set \(Z_\tau\) is disjoint from \(B\cup R_{\beta,J}\)
because \(Z_\tau\subseteq A_\tau\), the supports are contained in
\(P_\tau\), and \(P_\tau\cap A_\tau=\emptyset\).  Hence the full cylinder
factor splits into the removed factors and the residual block factor:
\[
p^{|P_\tau|}q^{|A_\tau|}
=
p^{u_R}p^{u_B}q^{|Z_\tau|}\,
  \nu(\iota(\beta,\tau)).
\]
The absent blue factors in \(A_\tau\) are contained in
\(\nu(\iota(\beta,\tau))\), as required by the normalization argument above.
The equality is interpreted with the zero convention already discussed when
\(p=0\) and support factors are present.
Substituting \(q=1-p\) gives the claimed inequality.  The complete terminal product law from
\(\noderef{FixedSetHistoryCellTerminalProductPartition}\) identifies these
path products with the terminal Bernoulli cylinder masses; the support
geometry comes from
\(\noderef{FixedSetHistoryCellRedBranchTranscriptCells}\); and the
functional disjointness used for injectivity is
\(\noderef{FixedSetHistoryCellRedFunctionalScanDisjointness}\).  Thus the
common fixed-\((B,J)\) residual leaf system is well-defined, normalized,
injective on realized cells, and satisfies the termwise domination.
\end{proof}
